\chapter{Osmolality (Serum and Urine)}

\section*{What Is This Test?}
The osmolality test measures the concentration of osmotically active particles (solutes) in serum, plasma, or urine. Osmolality is the number of osmoles (Osm) of solute per kilogram of solvent (water), expressed as mOsm/kg H\textsubscript{2}O. This is distinct from osmolarity, which is solute per liter of solution (mOsm/L)---the two are nearly identical in dilute solutions but diverge in conditions where plasma water content changes (hyperlipidemia, hyperproteinemia). Serum osmolality is tightly regulated by antidiuretic hormone (ADH/vasopressin) and thirst, maintaining a normal range of 275--295 mOsm/kg. The major contributors to serum osmolality are sodium (and its accompanying anions chloride and bicarbonate), accounting for approximately 92\% of serum osmolality, followed by glucose, and blood urea nitrogen (BUN). Calculated serum osmolality = 2[Na] + Glucose(mg/dL)/18 + BUN(mg/dL)/2.8. The osmolal gap is the difference between measured and calculated osmolality (normally <10 mOsm/kg). An elevated osmolal gap indicates the presence of an unmeasured osmole, most importantly toxic alcohols (methanol, ethylene glycol, isopropanol). Urine osmolality reflects the kidney's ability to concentrate or dilute urine and is essential for evaluating water homeostasis and polyuria/polydipsia syndromes.

\section*{Alternative Names}
Serum Osmolality; Plasma Osmolality; Urine Osmolality; Osmolal Gap; Osmolar Gap; Measured Osmolality; Calculated Osmolality

\section*{Principle}
Osmolality is measured by freezing point depression osmometry. The freezing point of pure water is 0.00 degrees C. The addition of 1 osmole of solute particles per kilogram of water depresses the freezing point by 1.86 degrees C. The instrument supercools the sample to approximately -7 degrees C and then induces crystallization by rapid agitation; the temperature rises to the freezing plateau, and the freezing point is precisely measured with a thermistor. The osmolality is calculated as: mOsm/kg = $\Delta$T / 1.86. Freezing point depression osmometry measures all osmotically active solutes regardless of their identity, including sodium, potassium, chloride, bicarbonate, glucose, urea, and unmeasured osmoles (ethanol, methanol, ethylene glycol, isopropanol, mannitol, sorbitol, glycine). This is the gold standard clinical method. Vapor pressure osmometry is an alternative but does not detect volatile solutes such as ethanol or methanol, limiting its clinical utility in toxicology. The osmolal gap = measured osmolality -- calculated osmolality. Calculated osmolality uses the equation: 2[Na(mEq/L)] + [Glucose(mg/dL)]/18 + [BUN(mg/dL)]/2.8. A more complete formula includes ethanol: 2[Na] + Glucose/18 + BUN/2.8 + Ethanol/4.6. This test is NOT performed by hematology analyzers.

\section*{History}
The concept of osmotic pressure was discovered by Jean-Antoine Nollet in 1748. Jacobus Henricus van 't Hoff received the first Nobel Prize in Chemistry (1901) for his work on osmotic pressure. Freezing point depression osmometry was developed in the 1950s. The clinical use of the osmolal gap for toxic alcohol detection was established in the 1980s. The relationship between plasma osmolality, ADH, and thirst was characterized by Andersson, Verney, and others in the 1950s-1960s.

\section*{Why This Test Is Ordered}
\begin{itemize}
  \item Evaluation of hyponatremia: determine whether hyponatremia is hypotonic (low osmolality), isotonic (normal osmolality---pseudohyponatremia from hyperlipidemia/hyperproteinemia), or hypertonic (high osmolality---hyperglycemia, mannitol)
  \item Detection of toxic alcohol ingestion: elevated osmolal gap (>10 mOsm/kg) in a patient with altered mental status, high-AG metabolic acidosis, or unexplained neurological symptoms suggests methanol, ethylene glycol, or isopropanol
  \item Diagnosis and classification of diabetes insipidus, SIADH, and primary polydipsia using simultaneous serum and urine osmolality measurements
  \item Evaluation of polyuria (>3 L/day urine output): differentiate water diuresis (primary polydipsia, diabetes insipidus) from solute diuresis (hyperglycemia, mannitol, high-protein tube feeds)
  \item Assessment of hypernatremia: hypernatremia always indicates hyperosmolality
  \item Monitoring of mannitol therapy for cerebral edema: target serum osmolality 300--320 mOsm/kg
  \item Water deprivation test for diabetes insipidus workup: serial serum and urine osmolality measurements following water restriction
\end{itemize}

\section*{Is This a Primary or Secondary Test}
Serum osmolality is a secondary test, ordered when initial electrolyte panels reveal hyponatremia, hypernatremia, or suspected toxic ingestion. It is not part of routine chemistry panels.

\section*{How This Test Is Performed}
For serum osmolality, a healthcare professional draws blood from a vein into a serum separator tube (gold-top) or lithium heparin tube (green-top). A tourniquet is briefly applied. For urine osmolality, a random (spot) urine sample is collected in a clean container. No preservatives are needed. The sample is analyzed by freezing point depression osmometry. Results are typically available within 1--2 hours. The tube must be capped to prevent evaporation, which would concentrate the sample and falsely elevate measured osmolality.

\section*{What Preparation Is Needed}
No fasting is required. However, avoid excessive water intake immediately before the test, as this can dilute serum osmolality. Inform the provider of recent ethanol ingestion (which contributes to measured osmolality and osmolal gap). For water deprivation test, the patient must be NPO for 8--12 hours under close supervision. Medications affecting osmolality: mannitol, IV contrast dye, sorbitol, lithium, diuretics, desmopressin. Urine osmolality is affected by fluid intake---a spot urine osmolality is most meaningful when interpreted alongside simultaneous serum osmolality and the clinical volume status.

\section*{Contraindications and Precautions}
No absolute contraindications. Important considerations: (1) The calculated osmolality equation assumes normal values for sodium, glucose, and BUN. Abnormal protein or lipid levels do not directly affect measured osmolality (they contribute minimally to osmolarity due to their large molecular size) but can cause pseudohyponatremia with indirect ISE, leading to a discrepancy between measured and calculated osmolality. (2) Ethanol is the most common cause of an elevated osmolal gap in emergency departments. Ethanol can be measured and included in the calculated osmolality: ethanol (mg/dL) / 4.6 = ethanol contribution in mOsm/kg. (3) The osmolal gap is a screening tool, not a definitive diagnostic test. Sensitivity for toxic alcohol detection is approximately 85--90\% when the gap exceeds 10 mOsm/kg. However, a normal osmolal gap does NOT exclude toxic alcohol ingestion, particularly if ingestion occurred many hours earlier (the parent alcohol has been metabolized to acids). (4) Isopropanol (rubbing alcohol) causes an elevated osmolal gap without metabolic acidosis (it is metabolized to acetone, not an acid). (5) Mannitol, sorbitol, glycine (irrigation fluid in TURP), and IVIG (maltose/sucrose stabilizers) can all increase the osmolal gap. (6) Ketoacidosis (DKA, alcoholic ketoacidosis) and lactic acidosis can mildly increase measured osmolality and the osmolal gap. (7) Chronic renal failure: uremic toxins increase measured osmolality slightly, so the osmolal gap may be mildly elevated. (8) A normal osmolal gap varies by laboratory: commonly <10 mOsm/kg, but some sources use <12--15 mOsm/kg. (9) Samples must not be exposed to air---evaporation concentrates the sample and raises measured osmolality.

\section*{Risks}
Minimal risk from venipuncture. Mild pain or bruising at the site resolves in 1--2 days. No risk for urine collection. Blood volume drawn (3--5 mL) is negligible.

\section*{What Happens After the Test}
Results are available within 1--2 hours. Interpretation depends on the clinical context: (1) Hyponatremia + low serum osmolality (<275 mOsm/kg): true hypotonic hyponatremia. Next, assess volume status and urine osmolality/urine Na. Urine osmolality <100 mOsm/kg suggests primary polydipsia; urine osmolality >100 mOsm/kg with UNa >20--40 suggests SIADH. (2) Hyponatremia + normal osmolality (275--295): pseudohyponatremia (hyperlipidemia, hyperproteinemia) or isotonic rinse solutions. (3) Hyponatremia + high osmolality (>295): hypertonic hyponatremia (hyperglycemia, mannitol). Corrected Na = measured Na + 1.6 $\times$ (glucose -- 100)/100. (4) Osmolal gap >10 mOsm/kg: consider toxic alcohols. Calculate contribution: assume 1 mOsm/kg per 1 mmol/L of toxic alcohol. Methanol (MW 32): serum level (mg/dL)/3.2 approximates contribution. Ethylene glycol (MW 62): serum level (mg/dL)/6.2. (5) Urine osmolality 300--900 mOsm/kg: normal random urine. Urine osmolality after 12-hour water deprivation should be >800 mOsm/kg (indicates intact renal concentrating ability). Urine osmolality <300 mOsm/kg with hypernatremia/hyperosmolality suggests diabetes insipidus.

\section*{Understanding Results}

\subsection*{Below Normal}
\begin{itemize}
  \item \textbf{Low serum osmolality (<275 mOsm/kg):} Indicates hypotonic hyponatremia; excess water relative to solute. Most common cause is SIADH.
  \item \textbf{SIADH:} Euvolemic hyponatremia; serum osmolality <275 mOsm/kg; urine osmolality inappropriately >100 mOsm/kg (often > serum osmolality); urine Na >20--40 mEq/L. Causes: CNS disorders, pulmonary disease, malignancy (especially small cell lung cancer), medications (SSRIs, carbamazepine, cyclophosphamide).
  \item \textbf{Primary polydipsia (psychogenic polydipsia, beer potomania):} Water intake exceeds 18--20 L/day; serum osmolality low; urine osmolality maximally dilute (<100 mOsm/kg, often <50 mOsm/kg); urine Na low (<20 mEq/L).
  \item \textbf{Cerebral salt wasting:} Hypovolemic hyponatremia with low serum osmolality; urine osmolality variable; urine Na >20--40 mEq/L. Seen with subarachnoid hemorrhage, neurosurgery, head trauma.
  \item \textbf{Hypothyroidism and adrenal insufficiency:} Euvolemic hyponatremia.
  \item \textbf{Postoperative state:} Transient SIADH from pain, nausea, narcotics; serum osmolality falls.
\end{itemize}

\subsection*{Above Normal}
\begin{itemize}
  \item \textbf{High serum osmolality (>295 mOsm/kg):} Indicates hypertonic state, almost always due to hypernatremia, hyperglycemia, or toxic solutes.
  \item \textbf{Hypernatremia (>145 mEq/L):} Osmolality approximately = 2 $\times$ Na + 10. Na 150 $\rightarrow$ osmolality ~310; Na 160 $\rightarrow$ osmolality ~330; Na 170 $\rightarrow$ osmolality ~350. Causes: water loss (GI, renal, insensible), inadequate water intake, diabetes insipidus, hypertonic saline administration.
  \item \textbf{Hyperglycemia:} DKA and HHS. Each 100 mg/dL increase in glucose above 100 mg/dL increases osmolality by ~5.5 mOsm/kg. In HHS, glucose often >600--1000 mg/dL, producing osmolality >320--350 mOsm/kg. Effective osmolality (tonicity) = measured osmolality -- BUN/2.8 (urea is an ineffective osmole, freely crossing cell membranes).
  \item \textbf{Elevated osmolal gap (>10 mOsm/kg):} Unmeasured osmoles. Methanol (MW 32---formic acid production, visual disturbances, retinal toxicity; osmolal gap appears within 1--2 hours of ingestion). Ethylene glycol (MW 62---glycolic and oxalic acid production, calcium oxalate crystals in urine, AKI). Isopropanol (MW 60---metabolized to acetone, CNS depression without acidosis; osmolal gap and ketosis without acidosis are classic). Ethanol (contributes ~22 mOsm/kg per 100 mg/dL ethanol). Mannitol (therapeutic for cerebral edema). Propylene glycol (vehicle for IV medications like lorazepam, diazepam; prolonged high-dose infusions). Glycine or sorbitol irrigating solutions (TURP syndrome). IVIG (maltose/sucrose stabilizers). Acetone (in DKA or isopropanol ingestion; contributes to both osmolal gap and ketosis).
  \item \textbf{Diabetes insipidus (DI):} Polyuria, hypernatremia, hyperosmolality, dilute urine. Central DI (no ADH production): pituitary/hypothalamic tumors, neurosurgery, head trauma, sarcoidosis, histiocytosis. Nephrogenic DI (renal resistance to ADH): lithium, hypercalcemia, hypokalemia, CKD, genetic (aquaporin-2 or V2 receptor mutation). Water deprivation test: after dehydration, central DI shows urine osmolality that rises with desmopressin (DDAVP) administration; nephrogenic DI shows no rise with desmopressin.
\end{itemize}

\section*{Normal Results}
\begin{itemize}
  \item \textbf{Serum osmolality (adults):} 275--295 mOsm/kg H\textsubscript{2}O
  \item \textbf{Serum osmolality (children):} 275--290 mOsm/kg
  \item \textbf{Urine osmolality (random):} 300--900 mOsm/kg (varies widely with hydration)
  \item \textbf{Urine osmolality (after 12-14h water deprivation):} >800 mOsm/kg (normal concentrating ability)
  \item \textbf{Urine osmolality (maximal dilution):} 50--100 mOsm/kg (normal diluting ability)
  \item \textbf{Osmolal gap:} <10 mOsm/kg (some sources <12 mOsm/kg)
  \item \textbf{Calculated serum osmolality:} 2[Na] + Glucose/18 + BUN/2.8 + Ethanol/4.6
  \item \textbf{Effective osmolality (tonicity):} 2[Na] + Glucose/18 (excludes urea)
\end{itemize}

\section*{Follow-up Tests}
\begin{itemize}
  \item \textbf{Simultaneous urine osmolality and urine sodium:} Essential for hyponatremia workup. Urine osmolality >100 mOsm/kg with UNa >40 $\rightarrow$ SIADH; urine osmolality <100 mOsm/kg with UNa <20 $\rightarrow$ primary polydipsia.
  \item \textbf{Serum electrolytes (BMP/CMP):} Including sodium, glucose, BUN for calculated osmolality and osmolal gap.
  \item \textbf{Serum ethanol level:} Most common cause of elevated osmolal gap. Include in calculated osmolality formula: ethanol (mg/dL)/4.6.
  \item \textbf{Toxic alcohol screen (methanol, ethylene glycol, isopropanol):} Urgently if osmolal gap >10 mOsm/kg in patient with altered mental status, visual complaints, or high-AG metabolic acidosis.
  \item \textbf{Arterial blood gas (ABG):} Assess pH and bicarbonate. High-AG metabolic acidosis with elevated osmolal gap strongly suggests methanol or ethylene glycol.
  \item \textbf{Water deprivation test:} Formal ADH response testing. After 8--12 hour water deprivation (or when body weight decreases by 3--5\% or urine osmolality plateaus), administer desmopressin. Measure hourly urine osmolality and serum osmolality.
  \item \textbf{Copeptin level:} C-terminal fragment of pro-vasopressin; emerging marker for ADH secretion. Elevated copeptin with hypernatremia suggests appropriate ADH response (dehydration); low copeptin with hypernatremia suggests diabetes insipidus.
  \item \textbf{Renal function (BUN, creatinine):} Assess for renal concentrating ability.
  \item \textbf{Serum calcium:} Hypercalcemia causes nephrogenic DI.
  \item \textbf{Brain MRI (pituitary protocol):} Evaluate for pituitary/hypothalamic pathology in central DI.
\end{itemize}

\section*{References}
\begin{enumerate}
  \item Burtis CA, Ashwood ER, Bruns DE. Tietz Textbook of Clinical Chemistry and Molecular Diagnostics. 7th ed. Elsevier; 2023.
  \item Wallach J. Interpretation of Diagnostic Tests. 10th ed. Wolters Kluwer; 2022.
  \item Tierney LM, Saint S, Drazen JM. CURRENT Medical Diagnosis and Treatment. McGraw-Hill; 2023.
  \item Fischbach FT, Dunning MB. A Manual of Laboratory and Diagnostic Tests. 10th ed. Wolters Kluwer; 2023.
  \item Verbalis JG, Goldsmith SR, Greenberg A, et al. Diagnosis, evaluation, and treatment of hyponatremia: expert panel recommendations. Am J Med. 2013;126(10 Suppl 1):S1--S42.
  \item Kraut JA, Kurtz I. Toxic alcohol ingestions: clinical features, diagnosis, and management. Clin J Am Soc Nephrol. 2008;3(1):208--225.
  \item Sterns RH. Disorders of plasma sodium---causes, consequences, and correction. N Engl J Med. 2015;372(1):55--65.
  \item Adrogu\'{e} HJ, Madias NE. Hypernatremia. N Engl J Med. 2000;342(20):1493--1499.
\end{enumerate}

\clearpage
